%==============================================================================
\section{Hồi quy đa thức (Polynomial Regression)}
\label{sec:polynomial-regression}
%==============================================================================

\begin{dinhnghia}[title={Hồi quy đa thức}]
Mô hình quan hệ giữa biến độc lập và phụ thuộc bằng \textbf{đa thức bậc $n$}, vượt qua giới hạn tuyến tính của hồi quy đơn/đa biến tuyến tính thuần.
\end{dinhnghia}

\begin{ynghia}[title={Khi nào dùng?}]
Dữ liệu \textbf{phi tuyến} không khớp đường thẳng; hồi quy tuyến tính bỏ sót pattern cong.
\end{ynghia}

\tsimg{05_polynomial_regression/img01.jpg}

\subsection{Phương trình}

\begin{dinhnghia}[title={Đa thức bậc $n$}]
\[
y = w_0 + w_1 x + w_2 x^2 + \cdots + w_n x^n + \epsilon
\]
Bậc 2 (parabol):
\[
y = w_0 + w_1 x + w_2 x^2 + \epsilon
\]
\end{dinhnghia}

\subsection{Cách hoạt động}

\begin{phuongphap}[title={Biến đổi feature}]
Tạo feature $x, x^2, \ldots, x^n$ từ $x$ gốc, rồi huấn luyện \texttt{LinearRegression} trên không gian mở rộng --- tương tự MLR nhưng các cột đa thức phụ thuộc cùng một $x$ đầu vào.
\end{phuongphap}

\subsection{Triển khai Python}

\begin{vidu}[title={ice\_cream\_selling\_data.csv}]
49 mẫu: \texttt{Temperature}, \texttt{IceCreamSales}. File: \path{data/ice_cream_selling_data.csv}.
\end{vidu}

\begin{vidu}[title={Bước 1 --- Nạp và trực quan}]
\begin{lstlisting}
import numpy as np
import pandas as pd
import matplotlib.pyplot as plt
from sklearn.preprocessing import PolynomialFeatures

data = pd.read_csv('data/ice_cream_selling_data.csv')
X = data.iloc[:, 0].values.reshape(-1, 1)
y = data.iloc[:, 1].values

plt.scatter(X, y, color="green")
plt.xlabel("Temperature (C)")
plt.ylabel("Ice Cream Sales (units)")
plt.title("Original Data")
plt.show()
\end{lstlisting}
\tsimg{05_polynomial_regression/img02.jpg}
\end{vidu}

\begin{vidu}[title={Đặc trưng đa thức bậc 2}]
\begin{lstlisting}
degree = 2
poly_features = PolynomialFeatures(degree=degree)
X_poly = poly_features.fit_transform(X)
\end{lstlisting}
\end{vidu}

\begin{vidu}[title={Bước 2 --- Huấn luyện}]
\begin{lstlisting}
from sklearn.linear_model import LinearRegression
lr_model = LinearRegression()
lr_model.fit(X_poly, y)
\end{lstlisting}
\end{vidu}

\begin{vidu}[title={Bước 3 --- Dự đoán}]
\begin{lstlisting}
y_pred = lr_model.predict(X_poly)
df = pd.DataFrame({'Actual Values': y, 'Predicted Values': y_pred})
print(df.head())
\end{lstlisting}
\end{vidu}

\begin{vidu}[title={Bước 4 --- $R^2$}]
\begin{lstlisting}
from sklearn.metrics import r2_score
print("R2:", r2_score(y, y_pred))
\end{lstlisting}
Bài gốc: $R^2 \approx 0{,}932$ --- khoảng 93\% biến thiên được giải thích.
\end{vidu}

\begin{vidu}[title={Bước 5 --- Vẽ đường cong hồi quy}]
\begin{lstlisting}
plt.scatter(X, y, color="green")
plt.plot(X, y_pred, color='red',
         label=f'Polynomial Regression (degree={degree})')
plt.xlabel("Temperature (C)")
plt.ylabel("Ice Cream Sales (units)")
plt.legend()
plt.show()
\end{lstlisting}
\tsimg{05_polynomial_regression/img03.jpg}
\end{vidu}

\begin{vidu}[title={Dự đoán nhiệt độ mới}]
\begin{lstlisting}
X_new = np.array([[1.9929]])
X_new_poly = poly_features.transform(X_new)
print(lr_model.predict(X_new_poly))
\end{lstlisting}
\end{vidu}

\begin{luuy}[title={Bậc đa thức}]
Bậc quá cao dễ overfitting; nên thử cross-validation hoặc regularization khi mở rộng thực tế.
\end{luuy}
